\chapter{Configuration Management}
\label{ch:configuration}

\section{Configuration Sources}

DAGonStar reads configuration from an INI file, conventionally named
\texttt{dagon.ini}; the sample file \texttt{dagon.ini.sample} documents every
key recognized by the current implementation. Configuration is loaded by
\texttt{dagon.config.read\_config()}, which returns a nested dictionary of
strings:

\begin{lstlisting}[caption={Reading configuration programmatically.}]
from dagon.config import read_config

cfg = read_config("dagon.ini")
scratch = cfg["batch"]["scratch_dir_base"]
\end{lstlisting}

A \texttt{Workflow} may also receive an in-memory configuration dictionary
directly, bypassing the INI file entirely --- useful for unit tests or
programmatic workflow construction, and used extensively by the LLM and Web
task lessons in Part~II:

\begin{lstlisting}[caption={In-memory configuration.}]
from dagon import Workflow

workflow = Workflow("Experiment", config={
    "batch": {"scratch_dir_base": "/tmp/", "run_base": "", "threads": "1"},
    "dagon_service": {"use": "False"},
    "slurm": {"partition": ""},
})
\end{lstlisting}

\section{Minimal Configuration for Local Batch Workflows}

\begin{lstlisting}[style=ini]
[dagon_service]
route = http://localhost:57000
use = False

[ftp_pub]
ip = localhost

[batch]
scratch_dir_base=/tmp/
run_base=
remove_dir=False
threads=1

[slurm]
partition=
\end{lstlisting}

The current implementation provides safe defaults for missing optional
\texttt{dagon\_service} and \texttt{ftp\_pub} values, but production
workflows should keep these sections explicit, for clarity and for the sake
of a future reader auditing the configuration.

\section{Section-by-Section Reference}

\subsection{\texttt{[dagon\_service]}}

\begin{lstlisting}[style=ini]
[dagon_service]
route = http://localhost:57000
use = False
\end{lstlisting}

When \texttt{use = True}, \texttt{Workflow} attempts to register itself with
the DAGon service through \texttt{dagon.api.API}; the service then mediates
workflow/task registration, status updates, and some transversal workflow
lookups (\Cref{sec:service-architecture}). When disabled, DAGonStar runs as a
purely local, in-process engine --- the recommended mode for introductory
examples and unit tests, and the mode used throughout this book's local
exercises.

\subsection{\texttt{[ftp\_pub]}}

\begin{lstlisting}[style=ini]
[ftp_pub]
ip = localhost
\end{lstlisting}

Used by transversal workflow handling when task data must be downloaded from
a host registered through the service. Current defaults: host
\texttt{localhost}, user \texttt{anonymous}, empty password.
\begin{cautionbox}
Do not store production passwords in the repository. Use local, ignored
configuration, environment variables, or a site-specific secret-management
process.
\end{cautionbox}

\subsection{\texttt{[dagon\_ip]}}

\begin{lstlisting}[style=ini]
[dagon_ip]
ip = localhost
\end{lstlisting}

Present in the sample configuration for service/node-related workflows; not
required for basic local examples.

\subsection{\texttt{[batch]}}

\begin{lstlisting}[style=ini]
[batch]
scratch_dir_base=/tmp/
run_base=
remove_dir=False
threads=1
\end{lstlisting}

\begin{description}[leftmargin=2em, style=nextline]
  \item[\texttt{scratch\_dir\_base}] base directory used to create task
    scratch directories (\Cref{sec:execution-lifecycle}).
  \item[\texttt{run\_base}] optional format string for grouping a workflow
    run under a timestamped subdirectory; supports strftime-style tokens and
    the literal token \texttt{\_\_MILLIS\_\_}.
  \item[\texttt{remove\_dir}] documented as a configuration intent; in the
    current implementation, task garbage collection is controlled primarily
    by task reference counts and task flags rather than this key alone.
  \item[\texttt{threads}] used by staging command generation for parallel
    movement modes.
\end{description}

\begin{lstlisting}[style=ini, caption={Grouping a run under a timestamped subdirectory.}]
[batch]
scratch_dir_base=/tmp/dagon
run_base=%Y%m%d-%H%M%S-__MILLIS__
threads=4
\end{lstlisting}

\subsection{\texttt{[slurm]}}

\begin{lstlisting}[style=ini]
[slurm]
partition=
\end{lstlisting}

The configured \texttt{partition} value is used when generating Slurm staging
commands. Individual \texttt{SLURM} task objects may also receive Slurm
parameters directly at construction time (\Cref{ch:backends}):

\begin{lstlisting}
task = DagonTask(
    TaskType.SLURM, "model", "./run_model.sh",
    partition="compute", ntasks=16, time="02:00:00",
)
\end{lstlisting}

\subsection{LLM Provider Sections}
\label{sec:config-llm}

\texttt{TaskType.LLM} reads an OpenAI-compatible provider configuration from
a section named \texttt{[llm.<name>]}; an arbitrary number of provider
sections may coexist:

\begin{lstlisting}[style=ini]
[llm.example]
endpoint=https://api.example.org
api_key_env=DAGON_LLM_EXAMPLE_API_KEY
model=example-chat-model
\end{lstlisting}

\texttt{api\_key\_env} names a runtime environment variable holding the
credential; \texttt{api\_key} is available only for a local, ignored
configuration file and must never be committed. \Cref{ch:llm-tasks}
treats the full request/response model for this task type.

\subsection{Web Task Authentication}

\texttt{TaskType.WEB} authentication is supplied through
\emph{environment-variable names} embedded in the task specification, such as
\texttt{token\_env}, \texttt{username\_env}, \texttt{password\_env}, or
\texttt{value\_env}. As with LLM provider keys, literal HTTP tokens,
passwords, or API keys must never be placed in configuration samples or
workflow source files (\Cref{ch:web-tasks}).

\subsection{Cloud Provider Sections}

\begin{lstlisting}[style=ini]
[ec2]
key=<aws_access_key_id>
secret=<aws_secret_access_key>
region=

[digitalocean]
key=

[gce]
key=
secret=
project=
\end{lstlisting}

Cloud tasks use Apache Libcloud through \texttt{dagon.cloud}. All values are
provider-specific and must be treated as secrets; the sample placeholders are
to be replaced only in local, non-committed configuration files. Install
cloud dependencies with:

\begin{lstlisting}[style=bash]
python -m pip install "dagonstar[cloud]"
\end{lstlisting}

\subsection{Logging Sections}

The sample configuration includes a standard Python
\texttt{logging.config.fileConfig} setup:

\begin{lstlisting}[style=ini]
[loggers]
keys=root

[handlers]
keys=stream_handler

[formatters]
keys=formatter

[logger_root]
level=DEBUG
handlers=stream_handler

[handler_stream_handler]
class=StreamHandler
level=DEBUG
formatter=formatter
args=(sys.stderr,)

[formatter_formatter]
format=%(asctime)s %(name)-12s %(levelname)-8s %(message)s
\end{lstlisting}

When a \texttt{Workflow} is constructed without an explicit \texttt{config}
dictionary, it calls \texttt{fileConfig(config\_file)} internally; therefore,
a file-based configuration \emph{must} include valid logging sections for the
engine to initialize correctly.

\section{Runtime Environment Variables}

SKYCDS and Globus credentials are intentionally not hard-coded anywhere in
configuration; the current implementation reads them exclusively from the
process environment:

\begin{lstlisting}[style=bash]
export DAGON_SKYCDS_CLIENT_TOKEN=...
export DAGON_SKYCDS_CATALOG_TOKEN=...
export DAGON_SKYCDS_API_TOKEN=...
export DAGON_SKYCDS_IP=...

export DAGON_GLOBUS_TRANSFER_TOKEN=...
\end{lstlisting}

The current Globus transfer setup primarily uses the OAuth flow implemented
in \texttt{GlobusManager.\_\_init\_\_}, which prompts the user to authorize a
native application interactively.

\section{Optional Integration Dependencies}

\begin{lstlisting}[style=bash]
python -m pip install "dagonstar[docker]"
python -m pip install "dagonstar[globus]"
python -m pip install "dagonstar[api]"
python -m pip install "dagonstar[all]"
\end{lstlisting}

\section{Configuration as a Layered Override Mechanism}

It is worth being explicit about the precedence relationship between the two
configuration-loading paths \Cref{ch:configuration} introduced --- the
file-based \texttt{config\_file} path and the in-memory \texttt{config}
dictionary --- because conflating them is a frequent source of confusion
when a workflow's behavior unexpectedly diverges from what a committed
\texttt{dagon.ini} appears to specify.

When a \texttt{Workflow} is constructed with an explicit \texttt{config}
dictionary, that dictionary is used \emph{in its entirety}, and
\texttt{config\_file} --- even if also supplied, or left at its default
value of \texttt{"dagon.ini"} --- is not consulted at all for any section.
This is not a partial override or a merge of the two sources; it is a
complete substitution. The practical consequence, already exploited
throughout the LLM, Web, and Native task lessons (\Cref{ch:llm-tasks,ch:web-tasks,ch:native-tasks}),
is that an in-memory \texttt{config} dictionary must itself be \emph{complete}
with respect to every section the workflow's tasks will need --- a runtime-
constructed \texttt{llm.local\_mock} section is of no help if the same
in-memory dictionary omits the \texttt{[batch]} section the workflow's
\texttt{BATCH} producer task requires, since there is no fallback to the
INI file's corresponding section once the in-memory path has been chosen.

\begin{cautionbox}
This all-or-nothing semantics also explains a subtlety in the logging
behavior noted earlier in this chapter: \texttt{fileConfig(config\_file)} is
invoked only in the file-based path. A workflow constructed with an explicit
\texttt{config} dictionary does not automatically configure Python's
logging subsystem; if a workflow author wants both an in-memory
configuration and properly initialized logging, the logging setup must be
performed separately, before workflow construction, using the standard
library's \texttt{logging} module directly.
\end{cautionbox}

\section{Auditing a Configuration File for Reproducibility}

\Cref{ch:configuration} closes, in its final section, with a list of
configuration rules for reproducible science; it is useful to turn that list
into an explicit auditing procedure a researcher can apply to an existing,
inherited \texttt{dagon.ini} before relying on it.

\begin{enumerate}[leftmargin=1.8em]
  \item \textbf{Locate every secret-bearing key.} Search the file for
    \texttt{key}, \texttt{secret}, \texttt{api\_key} (as opposed to
    \texttt{api\_key\_env}), and \texttt{password}; any non-empty,
    non-placeholder value found should be treated as a credential leak and
    rotated, not merely removed from the file.
  \item \textbf{Confirm the service is disabled unless genuinely needed.}
    Check \texttt{[dagon\_service] use}; for the overwhelming majority of
    single-researcher, single-machine workflows discussed in this book,
    this should read \texttt{False} (\Cref{sec:service-architecture}).
  \item \textbf{Verify the scratch base has adequate, durable storage.} As
    \Cref{ch:getting-started} cautioned, a default of \texttt{/tmp/} is
    convenient for short tutorial runs but unsuitable for any workflow
    relying on checkpoint reuse across sessions.
  \item \textbf{Cross-check every \texttt{[llm.<name>]} and Web-task
    \texttt{*\_env} reference against the deployment environment.} A
    provider section naming \texttt{api\_key\_env=DAGON\_RESEARCH\_LLM\_KEY}
    is only as reproducible as the documentation, external to the
    repository, recording that this environment variable must be set
    before the workflow can run --- confirm such documentation exists.
  \item \textbf{Confirm the logging sections are present and valid} if the
    workflow does not supply an in-memory \texttt{config} dictionary,
    per the previous section's discussion of \texttt{fileConfig} behavior.
\end{enumerate}

\begin{practicebox}
This auditing procedure is deliberately framed as something a researcher
does \emph{once}, when first inheriting or reviewing a workflow's
configuration, and then records as part of the workflow's own documentation
--- precisely the discipline \Cref{ch:configuration}'s closing section asks
for when it distinguishes operational secrets from reproducibility
metadata. An audit that is performed but never written down provides no
benefit to the next researcher who needs to run the same workflow.
\end{practicebox}

\section{Configuration Sections Not Yet Exercised in This Book}

For completeness, and because Appendix~C consolidates every section into a
single reference document, it is worth noting explicitly which
configuration sections this book's worked examples do \emph{not} exercise,
so that a reader does not mistake their absence from the lesson chapters for
an absence from the engine itself. The cloud provider sections
(\texttt{[ec2]}, \texttt{[digitalocean]}, \texttt{[gce]}) and the
\texttt{[dagon\_ip]} section fall into this category: \texttt{CLOUD} tasks
are described architecturally in \Cref{ch:architecture} and referenced in
the backend-selection guidance of \Cref{ch:backends}, but, consistent with
this book's local-first verification philosophy
(\Cref{ch:examples-catalog}), no worked, runnable cloud example is given,
because doing so responsibly would require a live, billable cloud account
that cannot be assumed of every reader. A researcher who needs to configure
a cloud provider section should treat the placeholders in Appendix~C as a
starting structural template, to be completed and verified against the
target provider's own current documentation and credential model rather
than against this book, since provider-specific authentication mechanisms
are, by their nature, outside the scope of what the DAGonStar documentation
itself specifies.

\section{Configuration Rules for Reproducible Science}

\begin{practicebox}
For scientific workflows, configuration is itself part of the scientific
record and should be versioned with care:
\begin{itemize}[leftmargin=1.4em]
  \item commit safe templates such as \texttt{dagon.ini.sample};
  \item never commit real secrets;
  \item record the configuration values that affect scientific
    reproducibility --- scratch layout, model parameters, container images,
    cloud instance types, and Slurm resource requests;
  \item keep operational secrets strictly separate from reproducibility
    metadata, e.g.\ by referencing environment-variable \emph{names} in
    versioned configuration while the variable \emph{values} live only in
    the runtime environment.
\end{itemize}
\end{practicebox}

The full, consolidated configuration reference is reproduced in Appendix~C
for use as a standing lookup table while authoring workflows.
